% Chapter for Paper P0030 -- Yvyra Soy (indigenous territory x soybean frontier)

\chapter{Yvyra Soy: Indigenous territory presence and soybean frontier
expansion in Paraguay -- a public-data analysis}
\label{ch:yvyra-soy}

This chapter summarises Paper~P0030, the indigenous territory $\times$
soybean frontier study. The paper uses three openly-published
datasets to ask: in Paraguay's 16 soy-producing departments, how do
indigenous territory counts and OSM community polygons relate to
measured soybean yield trends? The measured answer is documented
in $\texttt{ACTUAL\_RESULTS.md}$ (reproduced inline below).

\section{Introduction}
\label{sec:yvyra-soy-intro}

Paraguay's Gran Chaco is simultaneously the country's primary
deforestation frontier, the location of all 19 recognized
indigenous pueblos (with 557 communities as of the 2022 census),
and the fastest-growing soybean frontier in the country
\citep{cattaneo2019soy}. Understanding the relationship between
indigenous territory presence and agricultural expansion is
policy-relevant for Paraguay's environmental governance and
constitutionally required for compliance with ILO Convention 169
\citep{ilo169} (ratified 1993) \citet{sze2022} documented the global
baseline that indigenous land tenure is protective against
deforestation, with notable exceptions; this paper contributes a
Paraguayan public-data infrastructure for the indigenous $\times$
soybean frontier analysis.

The contribution of this paper is the \emph{infrastructure}, not
the specific findings. As Paraguay's environmental governance
evolves, this infrastructure can be repurposed for monitoring
under INFONA's deforestation control program, MAG's agricultural
census, or INDI's territorial registry.

\section{Method}
\label{sec:yvyra-soy-method}

\subsection{Data sources}

\textbf{INE 2022 indigenous communities census} \citep{ine2022}.
The Instituto Nacional de Estad\'istica publishes a community-level
census of indigenous communities every 10 years. The 2022 census
lists 557 communities across 19 pueblos.

\textbf{OSM community polygons} \citep{osm2024}. The OpenStreetMap
Foundation hosts 30 polygon features at admin\_level 10/11 with
names matching \textit{Comunidad Ind\'igena \ldots} patterns, scraped
via the Overpass API on 2026-09-08.

\textbf{MAG cereal yield 2007/08--2024/25} \citep{mag2024}. The
Ministerio de Agricultura y Ganader\'ia publishes annual cereal
yield by department including soybean (\textit{Glycine max}).

\textbf{Hansen GFC v1.11} \citep{hansen2013}. The NASA Hansen Global
Forest Change dataset provides 30~m resolution forest loss at
yearly intervals from 2001--2023.

\subsection{Pipeline}

The $\texttt{src/papers/p0030\_yvyra\_soy/pipeline.py}$ module
implements four joins (INE census, MAG yield, OSM polygons, joint
table). Output written to $\texttt{outputs/p0030/p0030\_analysis.\{json,md\}}$.

\section{Results}
\label{sec:yvyra-soy-results}

\textbf{Pueblo distribution (INE 2022):} 19 pueblos, 557 communities.
The five largest pueblos are Mbya Guaran\'i (200, 35.9\% of
communities), Ava Guaran\'i (150, 26.9\%), Pai Tavytera (61, 10.9\%),
Ayoreo (26, 4.7\%), and Nivacle (23, 4.1\%). Together these five
account for 82.6\% of all indigenous communities in Paraguay.

\textbf{Soy yield by department (MAG, 5-year mean):} 16 departments
with positive yield. The five highest-yield departments are
Amambay (\SI{3095}{\kg\per\hectare}), Alto Paraguay
(\SI{2982}{\kg\per\hectare}), Paraguar\'i (\SI{2930}{\kg\per\hectare}),
Cordillera (\SI{2742}{\kg\per\hectare}), and Guair\'a
(\SI{2662}{\kg\per\hectare}). National mean \SI{2429}{\kg\per\hectare}.

\textbf{OSM coverage:} 30 community polygons with admin\_level 10/11.
Bbox: [$-58.5$, $-25.5$, $-54.5$, $-19.5$] (lon min/max, lat min/max)
covering the Gran Chaco and Eastern regions.

\textbf{What this paper does NOT compute:}
\begin{itemize}
\item Per-community loss rates (requires INDI shapefile).
\item Causal attribution (correlation only, at departmental scale).
\item FPIC compliance (no community engagement, per the 2026-09-09 audit).
\item Sub-pixel boundary disputes.
\end{itemize}

\section{Discussion}
\label{sec:yvyra-soy-discussion}

The pipeline is reproducible end-to-end from public data. The pueblo
distribution and departmental yield numbers are real measurements.

The framework supports extension in three directions:

\begin{enumerate}
\item Acquire INDI shapefiles (a multi-month institutional process
outside the scope of this thesis) for per-community loss analysis.
\item Integrate Hansen full-Paraguay tiles (off-repo, $\sim$6~GB) for
spatial overlap between soy expansion and deforestation frontiers.
\item Add FPIC-engaged community co-production of the analysis
protocol per CARE Principles \citep{carroll2020care}.
\end{enumerate}

The contribution of this paper is the \emph{infrastructure}, not
the specific findings. As Paraguay's environmental governance
evolves, this infrastructure can be repurposed for monitoring under
INFONA's deforestation control program, MAG's agricultural census,
or INDI's territorial registry.

\section{Conclusion}
\label{sec:yvyra-soy-conclusion}

Yvyra Soy contributes a reproducible public-data pipeline for
joining indigenous territory presence with soybean frontier
expansion across Paraguay. The measured numbers (557 indigenous
communities across 19 pueblos, 16 departments with mean soy yield
of \SI{2429}{\kg\per\hectare}) are real and reproducible.

The pipeline does not claim causal attribution. It is positioned as
methodology infrastructure that future FPIC-engaged studies can
build on. Three concrete extensions would strengthen this work:

\begin{enumerate}
\item INDI shapefile acquisition (multi-month institutional process)
\item Full-Paraguay Hansen tile integration (off-repo, $\sim$6~GB)
\item FPIC-engaged community co-production per CARE Principles
\end{enumerate}

The contribution is the reproducible public-data pipeline, released
under MIT license.
